\SourceLeaf{162}
\SixthArticle{30}
E formula 38 sequitur
\[
\frac{\Pi(k,z+\omega)}{\Pi(k,z)}
=\frac{z+1}{z+1+\omega}\cdot\frac{z+2}{z+2+\omega}\cdot\frac{z+3}{z+3+\omega}
\mathinner{.\,.\,.\,.}\frac{z+k}{z+k+\omega}\cdot k^{\omega}
\]
unde sumtis logarithmis, in series infinitas evolutis, prodit
\[
\begin{aligned}
{}\text{[60]}\quad&\log\Pi(k,z+\omega)=\log\Pi(k,z)\\
&+\omega\left(\log k-\frac{1}{z+1}-\frac{1}{z+2}-\frac{1}{z+3}-\text{etc.}-\frac{1}{z+k}\right)\\
&+\tfrac12\omega\omega\left(\frac{1}{(z+1)^2}+\frac{1}{(z+2)^2}+\frac{1}{(z+3)^2}+\text{etc.}+\frac{1}{(z+k)^2}\right)\\
&-\tfrac13\omega^3\left(\frac{1}{(z+1)^3}+\frac{1}{(z+2)^3}+\frac{1}{(z+3)^3}+\text{etc.}+\frac{1}{(z+k)^3}\right)\\
&+\text{etc. in inf.}
\end{aligned}
\]
Series, hic in $\omega$ multiplicata, quae, si magis placet, ita etiam exhiberi potest,
\[
\begin{aligned}
&-\frac{1}{z+1}+\log 2-\frac{1}{z+2}+\log\tfrac32-\frac{1}{z+3}\\[4pt]
&+\log\tfrac43-\frac{1}{z+4}+\log\tfrac54-\text{etc.}
+\log\frac{k}{k-1}-\frac{1}{z+k}
\end{aligned}
\]
e terminorum multitudine finita constat, crescente autem $k$ in infinitum, ad limitem certum converget, qui novam functionum transscendentium speciem nobis sistit, in posterum per $\Psi z$ denotandam.

Designando porro summas serierum sequentium, in \textit{infinitum} extensarum,
\[
\begin{array}{l}
\dfrac{1}{(z+1)^2}+\dfrac{1}{(z+2)^2}+\dfrac{1}{(z+3)^2}+\text{etc.}\\[10pt]
\dfrac{1}{(z+1)^3}+\dfrac{1}{(z+2)^3}+\dfrac{1}{(z+3)^3}+\text{etc.}\\[10pt]
\dfrac{1}{(z+1)^4}+\dfrac{1}{(z+2)^4}+\dfrac{1}{(z+3)^4}+\text{etc.}\\[10pt]
\text{etc.}
\end{array}
\]
resp. per $P,Q,R$ etc. (pro quibus signa functionalia introducere minus necessarium videtur), habebimus
\[
\begin{aligned}
{}\text{[61]}\quad&\log\Pi(z+\omega)=\log\Pi z+\omega\Psi z\\[4pt]
&+\tfrac12\omega\omega P-\tfrac13\omega^3Q+\tfrac14\omega^4R-\text{etc.}
\end{aligned}
\]
Manifesto functio $\Psi z$ erit functio derivata prima functionis $\log\Pi z$, adeoque
\[
{}\text{[62]}\qquad \frac{\mathrm{d}\Pi z}{\mathrm{d}z}=\Pi z\cdot\Psi z
\]
Perinde erit $P=\dfrac{\mathrm{d}\Psi z}{\mathrm{d}z}$, $Q=-\dfrac{\mathrm{dd}\Psi z}{2\mathrm{d}z^2}$, $R=+\dfrac{\mathrm{d}^3\Psi z}{2.3\mathrm{d}z^3}$ etc.
